\section{Il Processo di Compilazione}
Der Quellcode wird durch verschiedene Phasen in Binärdaten umgewandelt:

\begin{enumerate}
    \item \textbf{Präprozessor:} Verarbeitet Einbindungen und Makros. (\texttt{\#define}, \texttt{\#if}, \dots)
    \item \textbf{Compiler:} Übersetzt die Sprache auf hoher Ebene in Assembly.
    \item \textbf{Assembler:} Konvertiert Assembly in relokierbaren Maschinencode.
    \item \textbf{Linker/Binder:} Verbindet Objektdateien und Bibliotheken.
    \item \textbf{Loader:} Konvertiert relokierbare Adressen in absolute und lädt sie in den Speicher.
\end{enumerate}

\subsection{Analisi del Compilatore}
\begin{itemize}
    \item \textbf{Lexikalische Analyse:} Entfernen von Leerzeichen und Gruppierung von Tokens.
    \item \textbf{Syntaktische/Semantische Analyse:} Prüfung der Struktur und der Datentypen.
    \item \textbf{Optimierung:} Verbesserung des Codes hinsichtlich Geschwindigkeit oder Speicherplatz.
    \item \textbf{Generierung:} Erzeugung des Zielcodes.
\end{itemize}

\section{Interfacciamento I/O}
Die Peripheriegeräte werden über spezifische Register in den Speicher abgebildet (\textit{Memory Mapped I/O}):
\begin{itemize}
    \item \textbf{Datenregister:} Enthält die zu verarbeitenden Daten.
    \item \textbf{Steuerregister:} Konfiguration des Peripheriegeräts.
    \item \textbf{Statusregister:} Zeigt den aktuellen Zustand des I/O (z. B. Busy, Ready) an.
\end{itemize}


